\documentclass[11pt,reqno]{amsart}
\usepackage[T1]{fontenc}
\usepackage{lmodern,microtype,mathtools,amssymb}
\usepackage[a4paper,margin=27mm]{geometry}
\usepackage[hidelinks]{hyperref}
\usepackage{enumitem}
\setlist[enumerate]{leftmargin=*,itemsep=2pt,topsep=4pt}
\numberwithin{equation}{section}
\newtheorem{theorem}{Theorem}[section]
\newtheorem{lemma}[theorem]{Lemma}
\newtheorem{proposition}[theorem]{Proposition}
\newtheorem{corollary}[theorem]{Corollary}
\theoremstyle{remark}
\newtheorem{remark}[theorem]{Remark}
\newcommand{\C}{\mathbb C}
\newcommand{\N}{\mathbb N}
\newcommand{\ee}{\mathrm e}
\newcommand{\dist}{\operatorname{dist}}
\newcommand{\Zer}{\mathcal Z}
\newcommand{\Aab}{A_{a,b}}
\allowdisplaybreaks[1]
\emergencystretch=1.5em
\title[Formalized sparse Fock construction]{A formally verified sparse Fock construction for Erd\H{o}s Problem 906}
\author{Zhijie He}
\date{September 17, 2026}
\keywords{Entire functions, successive derivatives, sparse power series, formal verification, zero localization}
\hypersetup{pdftitle={A formally verified sparse Fock construction for Erdos Problem 906},pdfauthor={Zhijie He}}

\begin{document}

\begin{abstract}
I study the explicit entire function
\[
 F(z)=\sum_{j\ge 1}\frac{z^{\lfloor j^{3/2}\rfloor}}{\sqrt{\lfloor j^{3/2}\rfloor!}}.
\]
Every nonempty open subset of the plane contains a zero of $F^{(n)}$ for every sufficiently large $n$.  On each fixed annulus the zeros admit a quantitative description: they lie in pairwise disjoint disks around explicit two-term model points, each such disk contains exactly one simple zero, and the zero set has covering radius $O(n^{-1/6})$.  The construction is deterministic.  Every theorem, proposition, and corollary stated in this note has a corresponding declaration in the accompanying Lean~4/Mathlib development.  No claim of being the first solution of Erd\H{o}s Problem~906 is made.
\end{abstract}
\maketitle

\section{Statement and context}\label{sec:intro}
For a nonzero entire function $g$, write $\Zer(g)$ for its zero set.  The cofinite zero-hitting property is
\begin{equation}\label{eq:cofinite}
 \forall U\subset\C\text{ nonempty open},\quad
 \exists N(U)\quad\forall n\ge N(U):\quad \Zer(g^{(n)})\cap U\ne\varnothing.
\end{equation}
For a transcendental entire function this is equivalent to the formulation of Erd\H{o}s Problem~906 using every strictly increasing sequence of derivative orders.  Indeed, if \eqref{eq:cofinite} fails for an open set $U$, the infinitely many exceptional orders contain a strictly increasing sequence whose derivative-zero union misses $U$.  Conversely, a cofinite set of orders meets every infinite increasing sequence eventually.

Erd\H{o}s recorded the problem in \cite[p.~72, IV.1(i)]{Erdos82}.  Barth and Schneider \cite[Theorem~1]{BS72} solved a different interpolation problem: for prescribed discrete sets $S_k$ one may choose derivative orders $n_k$ so that $f^{(n_k)}$ vanishes on $S_k$.  That result does not contain the cofinite quantifier in \eqref{eq:cofinite}.  Gethner's final-set framework likewise concerns points met by zeros of infinitely many derivatives, a weaker quantifier than \eqref{eq:cofinite} \cite{Gethner85}.  Hou recently gave a probabilistic bounded-coefficient Fock-series construction satisfying \eqref{eq:cofinite}, together with a Lean formalization of his main theorem and growth bound \cite{Hou26}.  The point here is different: the function below is explicit and sparse, and the derivative zeros can be localized quantitatively.  I make no historical-priority claim beyond these stated features.

Put
\begin{equation}\label{eq:series}
 \nu_j=\lfloor j^{3/2}\rfloor,\qquad
 q_j=\nu_{j+1}-\nu_j,\qquad
 F(z)=\sum_{j=1}^{\infty}\frac{z^{\nu_j}}{\sqrt{\nu_j!}}.
\end{equation}
For $0<a\le b$, set
\[
 \Aab=\{z\in\C:a\le |z|\le b\}.
\]
For an active index $j$ with $\nu_j\ge n$, define
\begin{equation}\label{eq:weights}
 m_j=\nu_j-n,\qquad
 A_j=\frac{\sqrt{\nu_j!}}{m_j!},\qquad
 T_j(z)=A_jz^{m_j},\qquad
 \rho_j=\left(\frac{A_j}{A_{j+1}}\right)^{1/q_j}.
\end{equation}
The radius $\rho_j$ is where $|T_j|=|T_{j+1}|$.  The corresponding nonzero zeros of the two-term model $T_j+T_{j+1}$ are
\begin{equation}\label{eq:centers}
 z_{j,\ell}=\rho_j\exp\!\left(\frac{(2\ell+1)\pi i}{q_j}\right),
 \qquad 0\le \ell<q_j.
\end{equation}
For $c>0$ write
\begin{equation}\label{eq:modeldisk}
 \mathcal D_{j,\ell}(c)=\overline D\!\left(z_{j,\ell},\frac{c\rho_j}{q_j}\right).
\end{equation}

\begin{theorem}[Cofinite derivative zeros]\label{thm:cofinite}
For every nonempty open set $U\subset\C$ there is $N$ such that, for every $n\ge N$, the derivative $F^{(n)}$ has a zero in $U$.
\end{theorem}

\begin{theorem}[Quantitative annular covering]\label{thm:covering}
Let $0<a\le b$.  There are constants $C>0$ and $N$ such that for every $n\ge N$ and every $w\in\Aab$ there is a zero $z$ of $F^{(n)}$ satisfying
\begin{equation}\label{eq:coveringrate}
 |z-w|\le C n^{-1/6}.
\end{equation}
\end{theorem}

\begin{proposition}[One simple zero in each model disk]\label{prop:modeldisk}
Let $0<a\le b$.  There exist $c\in(0,1]$ and $N$ such that for every $n\ge N$, every active $j$ with $a\le\rho_j\le b$, and every $0\le\ell<q_j$, the disk $\mathcal D_{j,\ell}(c)$ contains exactly one zero of $F^{(n)}$.  That zero is simple.
\end{proposition}

\begin{proposition}[Exhaustion of annular zeros]\label{prop:exclusion}
Let $0<a\le b$ and $c>0$.  For all sufficiently large $n$, every zero $z$ of $F^{(n)}$ with $a\le |z|\le b$ belongs to a disk $\mathcal D_{j,\ell}(c)$ with
\[
 \frac a2\le\rho_j\le 2b,
 \qquad 0\le\ell<q_j.
\]
\end{proposition}

\begin{corollary}[Eventual annular simplicity]\label{cor:simplicity}
Let $0<a\le b$.  For all sufficiently large $n$, every zero of $F^{(n)}$ in $\Aab$ is simple.
\end{corollary}

\begin{proposition}[Pairwise disjoint model disks]\label{prop:disjoint}
Let $0<a\le b$.  If $0<c\le1$, then for all sufficiently large $n$ the admissible disks $\mathcal D_{j,\ell}(c)$ with $a\le\rho_j\le b$ are pairwise disjoint.
\end{proposition}

\begin{proposition}[Finite model-disk zero count]\label{prop:count}
Let $0<a\le b$.  There exist $c\in(0,1]$ and $N$ such that, for every $n\ge N$ and every finite set $I$ of admissible pairs $(j,\ell)$ with $a\le\rho_j\le b$, the union of the disks indexed by $I$ contains exactly $|I|$ distinct zeros of $F^{(n)}$, all simple.
\end{proposition}

The remainder of the paper proves these statements in a form that mirrors the structure of the formal development.  The exact order and type, sectorial zero asymptotics, limiting zero measures, other values of the support exponent, and the critical endpoint $p=4/3$ are deliberately not part of this submission.

\section{Sparse support and derivative weights}\label{sec:weights}
The two relevant scales are immediate from the support.  The mean-value theorem and the floor error give
\begin{equation}\label{eq:gap}
 \frac32\sqrt j-1\le q_j\le \frac32(\sqrt j+1)+1\qquad(j\ge1).
\end{equation}
Thus, when $\nu_j=n+O(\sqrt n)$,
\begin{equation}\label{eq:qscale}
 j\asymp n^{2/3},\qquad q_j\asymp n^{1/3}.
\end{equation}
In particular the first active exponent $m_{j_0}$, where $j_0=\min\{j:\nu_j\ge n\}$, is $O(n^{1/3})$.

The series in \eqref{eq:series} is entire.  A convenient uniform majorant is
\begin{equation}\label{eq:growthupper}
 |F(z)|\le \sqrt{1+|z|^2}\,\exp\!\left(\frac{1+|z|^2}{2}\right).
\end{equation}
Indeed, the sparse series is termwise dominated by $\sum_{k\ge0}|z|^k/\sqrt{k!}$, and a Cauchy--Schwarz estimate gives \eqref{eq:growthupper}.  Termwise differentiation is therefore legitimate, and
\begin{equation}\label{eq:derivative}
 F^{(n)}(z)=\sum_{\nu_j\ge n} A_jz^{m_j}.
\end{equation}
The support is infinite, hence $F$ is not a polynomial.  In the Lean development transcendence is certified directly by showing that sufficiently high derivatives of a polynomial vanish while the corresponding derivatives of $F$ do not.

For integers $k\ge n$ and $r>0$, set
\begin{equation}\label{eq:Phi}
 \Phi_{n,r}(k)=\frac12\log(k!)-\log((k-n)!)+(k-n)\log r.
\end{equation}
Its first difference is
\begin{equation}\label{eq:difference}
 d_{n,r}(k)=\Phi_{n,r}(k+1)-\Phi_{n,r}(k)
 =\log r+\frac12\log(k+1)-\log(k+1-n).
\end{equation}
Writing $m=k-n$,
\begin{equation}\label{eq:curvature}
 d_{n,r}(k)-d_{n,r}(k+1)
 =\log\!\left(1+\frac1{m+1}\right)
 -\frac12\log\!\left(1+\frac1{n+m+1}\right)>0.
\end{equation}
Hence $\Phi_{n,r}$ is strictly discretely concave.  Uniformly for $m=O(\sqrt n)$, the drop in \eqref{eq:curvature} is comparable to $n^{-1/2}$.

The crossing radius has the exact representation
\begin{equation}\label{eq:rhoexact}
 \log\rho_j=\frac1{q_j}\sum_{s=1}^{q_j}
 \left\{\log(m_j+s)-\frac12\log(n+m_j+s)\right\}.
\end{equation}
The summand is increasing in $s$, so the $\rho_j$ are strictly increasing.  In the saddle range $m_j=O(\sqrt n)$, \eqref{eq:rhoexact} and \eqref{eq:qscale} give
\begin{equation}\label{eq:rhoapprox}
 \rho_j=\frac{m_j}{\sqrt n}+O(n^{-1/6}),
\end{equation}
and on every fixed compact annulus
\begin{equation}\label{eq:rhospace}
 c_{a,b}n^{-1/6}\le \rho_{j+1}-\rho_j\le C_{a,b}n^{-1/6}
\end{equation}
whenever both radii lie in a fixed enlargement of the annulus.  These estimates explain the exponent in Theorem~\ref{thm:covering}: the radial mesh is $n^{-1/6}$, while the angular mesh on a crossing circle is $q_j^{-1}=O(n^{-1/3})$.

\section{Two-term localization}\label{sec:twoterm}
Fix a compact saddle range $m_j\le B\sqrt n$.  At $r=\rho_j$ the $q_j$ consecutive differences
\[
 d_{n,r}(\nu_j),\ldots,d_{n,r}(\nu_j+q_j-1)
\]
have sum zero.  By \eqref{eq:curvature}, their endpoint values have opposite signs and magnitude $\gg q_j/\sqrt n$.  If
\[
 |q_j\log(|z|/\rho_j)|\le L,
\]
the same signs persist for large $n$.  Since every further sparse step has length $\gg n^{1/3}$, one gains a logarithmic loss
\begin{equation}\label{eq:En}
 E_n\asymp \frac{q_j^2}{\sqrt n}\asymp n^{1/6}
\end{equation}
on each step away from the pair $j,j+1$.  Summing the resulting geometric tails yields
\begin{equation}\label{eq:tail}
 \frac{\sum_{h\ne j,j+1}|T_h(z)|}{|T_j(z)|}
 \le (1+\ee^L)\frac{\ee^{-E_n}}{1-\ee^{-E_n}}.
\end{equation}
This is the quantitative reduction to two adjacent terms.

Let $z_0=z_{j,\ell}$ and $q=q_j$.  In the logarithmic coordinate $z=z_0\ee^w$, division by the nonvanishing principal monomial puts the derivative in the form
\begin{equation}\label{eq:logmodel}
 \Psi(w)=1-\ee^{qw}+E(w).
\end{equation}
Fix a small $c_0>0$ and put $\zeta=c_0/q$.  The tail estimate \eqref{eq:tail}, applied on $|w|\le\zeta$, gives $|E(w)|\le K_n$ with $K_n=O(\ee^{-E_n})$.  At the centre $|\Psi(0)|\le K_n$.  On $|w|=\zeta$, the elementary estimate for the exponential remainder gives
\[
 |1-\ee^{qw}|\ge c_0/2
\]
when $c_0\le1/2$.  Hence $|\Psi|\ge c_0/2-K_n>K_n$ on the boundary for large $n$.  If $\Psi$ had no zero in $|w|\le\zeta$, the maximum-modulus principle applied to $1/\Psi$ would force the modulus at the centre to be at least the boundary minimum, a contradiction.  Thus a true zero lies within $O(\rho_j/q_j)$ of $z_0$.

It remains to prove uniqueness and simplicity without a zero-counting black box.  Write $t=z/z_0$ and normalize once more:
\begin{equation}\label{eq:tmodel}
 G(t)=1-t^q+R(t).
\end{equation}
For a fixed sufficiently small $s\asymp q^{-1}$, the same tail estimate holds on $|t-1|\le2s$ with $|R(t)|\le K_n$.  Cauchy's estimate therefore gives
\begin{equation}\label{eq:Rprime}
 |R'(t)|\le K_n/s\qquad(|t-1|\le s).
\end{equation}
Choose $s$ so that $qs\le1/8$.  Then $t^{q-1}$ stays within $1/7$ of $1$ on this disk.  Consequently
\begin{equation}\label{eq:powdiff}
 |t_1^q-t_2^q|\ge \frac{6q}{7}|t_1-t_2|
 \qquad(|t_i-1|\le s),
\end{equation}
while \eqref{eq:Rprime} gives
\[
 |R(t_1)-R(t_2)|\le (K_n/s)|t_1-t_2|.
\]
For large $n$, $K_n<qs/2$.  If $G(t_1)=G(t_2)=0$, comparison with \eqref{eq:powdiff} forces $t_1=t_2$.  Moreover
\[
 |q t^{q-1}|\ge 6q/7>K_n/s,
\]
so $G'(t)\ne0$ at its zero.  The zero is therefore unique and simple.  This proves Proposition~\ref{prop:modeldisk} by exactly the existence--uniqueness mechanism used in the formal development.

\section{Exhaustion, disjointness, and covering}\label{sec:covering}
Fix $r$ in a compact annulus and choose a summand of maximal modulus on $|z|=r$.  The normalized block slopes satisfy
\begin{equation}\label{eq:blockslope}
 \frac1{q_j}\log\frac{|T_{j+1}(z)|}{|T_j(z)|}
 =\log r-\log\rho_j.
\end{equation}
Because $\rho_j$ is strictly increasing, these slopes strictly decrease with $j$.  Hence the sparse terms increase to a single maximal region and then decrease.

Suppose $r$ lies between two consecutive crossing radii.  If $r$ is outside fixed logarithmic neighborhoods of both crossings, the middle term dominates both neighbors; the tail estimate then makes it larger than the sum of all remaining terms.  No zero can occur there.  If $r$ lies in one of the transition neighborhoods, the two-term estimate applies.  At a zero, division by the dominant monomial and \eqref{eq:tmodel} force the phase to lie within $O(1/q_j)$ of an odd $q_j$-th root of $-1$, hence the zero lies in a model disk.  Since \eqref{eq:rhoapprox} keeps the relevant crossing radius in $[a/2,2b]$, this proves Proposition~\ref{prop:exclusion}.

Apply Proposition~\ref{prop:modeldisk} on the enlarged annulus $A_{a/2,2b}$ and choose its disk constant in Proposition~\ref{prop:exclusion}.  Every annular zero then lies in a model disk whose unique zero is simple.  This proves Corollary~\ref{cor:simplicity}.

For the disjointness assertion, model centers on one circle satisfy
\begin{equation}\label{eq:angularsep}
 |z_{j,\ell}-z_{j,\ell'}|\ge \frac{4\rho_j}{q_j}
 \qquad(\ell\ne\ell'),
\end{equation}
for all sufficiently large relevant $j$.  With $c\le1$, disks on the same circle are therefore disjoint.  Disks on different circles are separated by the radial gap \eqref{eq:rhospace}, of order $n^{-1/6}$, whereas their radii are $O(q_j^{-1})=O(n^{-1/3})$.  Hence the whole admissible family is eventually pairwise disjoint, proving Proposition~\ref{prop:disjoint}.  Choosing the disk constant supplied by Proposition~\ref{prop:modeldisk}, disjointness and uniqueness give one distinct simple zero per disk in any finite admissible family.  This proves Proposition~\ref{prop:count}.

It remains to prove Theorem~\ref{thm:covering}.  Let $w\in\Aab$ and put $r=|w|$.  Choose $j$ so that
\[
 \nu_j\le n+r\sqrt n<\nu_{j+1},
\]
with the first active index used in the harmless initial case.  Equations \eqref{eq:qscale} and \eqref{eq:rhoapprox} give
\[
 |\rho_j-r|\le Cn^{-1/6}.
\]
Choose the nearest model angle on the circle $|z|=\rho_j$.  Its angular displacement from $w$ is $O(q_j^{-1})=O(n^{-1/3})$.  Apply Proposition~\ref{prop:modeldisk} on a fixed enlarged annulus containing all such crossing radii; it supplies a true zero within another $O(n^{-1/3})$.  The radial error dominates, giving \eqref{eq:coveringrate}.

Finally let $U$ be a nonempty open set.  It contains a nonzero point $w$ and a small closed disk about $w$ contained in some annulus $\Aab$.  For all sufficiently large $n$, Theorem~\ref{thm:covering} places a zero of $F^{(n)}$ in that disk.  This proves Theorem~\ref{thm:cofinite}.

\begin{remark}[The origin]
The simplicity statement is deliberately annular.  If $j_0=\min\{j:\nu_j\ge n\}$, then the multiplicity of the zero at the origin is exactly $\nu_{j_0}-n$.  For suitable derivative orders this multiplicity is large, so no global eventual-simplicity statement is true.
\end{remark}

\section{Formal verification and scope}\label{sec:formal}
The Lean development is pinned to Lean/Mathlib~4.28.0.  The main declarations corresponding to the statements above are:
\begin{center}
\begin{tabular}{p{0.42\textwidth}p{0.48\textwidth}}
\hline
Printed statement & Lean declaration \\
\hline
Entire function and derivative formula & \texttt{F\_differentiable}, \texttt{iteratedDeriv\_F} \\
Transcendence & \texttt{F\_not\_polynomial} \\
Upper growth bound \eqref{eq:growthupper} & \texttt{F\_norm\_le} \\
Theorem~\ref{thm:covering} & \texttt{annular\_covering\_rate\_p32} \\
Theorem~\ref{thm:cofinite} & \texttt{erdos906\_sparse\_fock\_p32} \\
Proposition~\ref{prop:modeldisk} & \texttt{model\_disk\_zero\_count\_one\_p32} \\
Proposition~\ref{prop:exclusion} & \texttt{annular\_zero\_exclusion\_p32} \\
Corollary~\ref{cor:simplicity} & \texttt{annular\_zeros\_simple\_deriv\_p32} \\
Proposition~\ref{prop:disjoint} & \texttt{model\_disks\_pairwise\_disjoint\_p32} \\
Proposition~\ref{prop:count} & \texttt{zero\_count\_model\_disk\_union\_p32} \\
\hline
\end{tabular}
\end{center}
The finite-union theorem uses ordinary cardinality of the set of distinct zeros; simplicity then makes this equal to total multiplicity on those disks.  The published source contains no \texttt{sorry}, \texttt{admit}, custom \texttt{axiom}, \texttt{unsafe}, \texttt{native\_decide}, or \texttt{@[implemented\_by]} in the project modules.  The main declarations are audited with \texttt{\#print axioms}; only the standard dependencies \texttt{propext}, \texttt{Classical.choice}, and \texttt{Quot.sound} occur.

The present note is intentionally narrower than the accompanying research directory.  Results there concerning general exponents, the endpoint $p=4/3$, sectorial counting, limiting zero measures, and support-profile limits are not claims of this formalized submission unless and until they receive separate machine verification.

\section*{AI assistance disclosure}
Language-model tools were used substantially during derivation, checking, exposition, and formalization.  The mathematical claims in this submission are separated from historical-priority claims, and the repository records the boundary of the machine-checked development.

\begin{thebibliography}{9}

\bibitem{BS72}
K.~F. Barth and W.~J. Schneider,
\emph{On a problem of Erd\H{o}s concerning the zeros of the derivatives of an entire function},
Proc. Amer. Math. Soc. \textbf{32} (1972), 229--232.
DOI: \href{https://doi.org/10.2307/2038336}{10.2307/2038336}.

\bibitem{Erdos82}
P. Erd\H{o}s,
\emph{Some of my favourite problems which recently have been solved},
in: L.~H.~Y. Chen, T.~B. Ng, M.~J. Wicks (eds.),
Proceedings of the International Mathematical Conference, Singapore 1981,
North-Holland Math. Stud. \textbf{74} (1982), 59--79.
DOI: \href{https://doi.org/10.1016/S0304-0208(08)70415-8}{10.1016/S0304-0208(08)70415-8}.

\bibitem{Gethner85}
R.~M. Gethner,
\emph{On the zeros of the derivatives of some entire functions of finite order},
Proc. Edinburgh Math. Soc. \textbf{28} (1985), 381--407.
DOI: \href{https://doi.org/10.1017/S001309150001720X}{10.1017/S001309150001720X}.

\bibitem{Hou26}
E. Hou,
\emph{Cofinite zeros of high derivatives},
arXiv:2607.20816 (2026).
\href{https://arxiv.org/abs/2607.20816}{arXiv:2607.20816}.

\end{thebibliography}

\end{document}
